\documentclass{article}
\usepackage{template/iclr2027_conference}
\usepackage{times}
\usepackage[T1]{fontenc}
\usepackage{amsmath,amssymb,graphicx,booktabs,longtable,array,calc}
\usepackage{hyperref}
\usepackage{etoolbox}
\hypersetup{colorlinks=true,linkcolor=blue,citecolor=blue,urlcolor=blue}
\providecommand{\tightlist}{\setlength{\itemsep}{0pt}\setlength{\parskip}{0pt}}
\makeatletter
\patchcmd{\@maketitle}{Under review as a conference paper at ICLR 2027}
  {Research draft --- not submitted}{\relax}{\errmessage{Draft header patch failed}}
\patchcmd{\@maketitle}{Anonymous authors\\Paper under double-blind review}
  {Research draft\\11 September 2026; not submitted for review}
  {\relax}{\errmessage{Draft author patch failed}}
\makeatother

\title{When Predictive Features Miss Motion: Measuring Temporal Accessibility in Skeleton JEPA}
\begin{document}
\maketitle
\begin{center}\textbf{Extended abstract}\end{center}
Predictive representations should be evaluated against observables with explicit meaning. We connect two gait studies and develop a measurement path toward selective future distillation.

In a retained skeleton-JEPA study with 625 clips from 93 source videos, trained predictors favor correct-clip hidden features over mismatched targets in all 375 diagnostic checks. Yet a common expanded readout of signed left--right coordinate-speed contrast falls from 0.222544 \(R^2\) at initialization to 0.100777--0.114209 across five training conditions. Every condition loses in every seed. Separately, a repaired RGB-conditioned experiment on 50 clips from 43 sources finds a skeleton increment of -0.00024242 \(R^2\), with conditional 95\% source-bootstrap interval {[}-0.00147402, +0.00091980{]}. The studies use different targets and \(R^2\) references. Neither establishes whether a skeleton-only student benefits from future video supervision.

We implement a source-held panel comparing a skeleton history block with current posture and observation support. Jointly regularized linear models, training-only preprocessing, matched temporal controls and exact baseline fallback constrain the comparison. The named history-block increment is +0.001994 \(R^2\), interval {[}-0.006673, +0.012608{]}, with 68.35\% positive paired draws: no supported temporal lead. This block also changes valid-conditioned confidence features and their regularization, so the estimate does not isolate motion from observation quality. The operational finite-model comparison does not estimate conditional mutual information.

Reflection and time reversal clarify the required observables. The existing laterality target changes sign under anatomical reflection but is unchanged by time reversal. Exact parity projections can enforce these rules even on zero features. A controlled 48-source fixture therefore tests utility separately: current posture and order-even summaries give \(R^2\) = 0 for a signed-velocity target, while ordered velocities recover it nearly exactly. This is a deliberately constructed readout calibration, not JEPA training or human-motion forecasting.

We propose selecting a small set of teacher directions by held-source temporal value, retaining an exact no-transfer option, then testing their benefit to a matched S-JEPA student. Future privileged distillation and equivariant features already have close precedents, notably SGDD and seq-JEPA. The prospective contribution is evidence that a student-aligned temporal selection criterion improves transfer or reliably rejects unsuitable supervision. That claim requires a complete real student comparison, common observable outcomes and independent source or participant confirmation. The current work preserves the negative findings as empirical motivation.

\section*{Draft status and AI use}
The proposed distillation method has no completed real-data student result. AI assistance supported evidence organization, literature checking, code and manuscript drafting. Human author verification remains required before submission.
\end{document}
